\notepage{74}{Degenerate perturbation theory}
When $E_n^0=E_m^0$, the criterion
\[
 \abs{\frac{\mel{m^0}{H'}{n^0}}{E_n^0-E_m^0}}\ll1
\]
breaks down. To understand the choice of basis, introduce
$H(\lambda)=H^0+\lambda H'$ and increase $\lambda$ from $0$ to $1$.
At $\lambda=0$ there is freedom to choose any orthonormal basis of a
degenerate eigenspace. Not every such basis connects smoothly to the
eigenstates selected when the perturbation is turned on.

\begin{center}
\begin{tikzpicture}[>=Latex,node distance=10mm]
 \node (a) {$\lambda=0:\quad\ket{\bar n^0},\ket{\bar m^0},\quad E_n^0=E_m^0$};
 \node[below=of a] (b) {$\lambda=1:\quad\ket n,\ket m,\quad E_n\ne E_m$};
 \draw[->] (a) -- node[right,font=\small] {suitable basis: smooth branches} (b);
\end{tikzpicture}
\end{center}
\smallheading{A two-state example}
Let
\[
 H^0=\begin{pmatrix}E&0\\0&E\end{pmatrix},\qquad
 H'=\begin{pmatrix}0&\Delta\\\Delta&0\end{pmatrix},\qquad\Delta\in\mathbb R.
\]
The exact eigenvalues follow from
\[
 \det\begin{pmatrix}E-\varepsilon&\Delta\\\Delta&E-\varepsilon\end{pmatrix}=0,
 \qquad\varepsilon=E\pm\Delta,
\]
with eigenstates
\[
 \ket{\pm}=\frac1{\sqrt2}\bigl(\ket{1^0}\pm\ket{2^0}\bigr).
\]
Choose instead the unperturbed basis
\[
 \ket{\bar1^0}=\frac{\ket{1^0}-\ket{2^0}}{\sqrt2},\qquad
 \ket{\bar2^0}=\frac{\ket{1^0}+\ket{2^0}}{\sqrt2}.
\]
In this ordered basis,
\[
 H'=\begin{pmatrix}-\Delta&0\\0&\Delta\end{pmatrix}.
\]
The perturbed eigenvectors are the same symmetric and antisymmetric
combinations, now already selected by the unperturbed basis.
\editorial{The source writes these basis vectors in the order
antisymmetric, symmetric, but reverses their eigenvalues in its drawn
diagonal matrix. The ordering is made consistent above.}
